\chapter{पृष्ठ}\label{ch:b2:surfaces}

वक्रों के बाद पृष्ठ: $\R^3$ में दो प्राचलों वाली वस्तुएँ।
\cref{ch:b2:diffcalc} का अवकल कलन हमें वह सब दे देता है जो चाहिए ---
आंशिक अवकलज स्पर्शी सदिश देते हैं, सदिश गुणनफल अभिलंब देता है,
\omterm{def:b2:linalg:det}{सारणिक} \omterm{def:b2:surfaces:area}{क्षेत्रफल} देते हैं। हम नियमित \omterm{def:b2:surfaces:param}{प्राचलित पृष्ठ}, उनके \omterm{def:b2:surfaces:tangent}{स्पर्श
समतल}, और \emph{\omterm{def:b2:surfaces:fff}{प्रथम आधारभूत रूप}} परिभाषित करते हैं, जो पृष्ठ पर
\omterm{def:b2:curves:length}{लंबाई} तथा \omterm{def:b2:surfaces:area}{क्षेत्रफल} के सारे मापन समेट लेता है। पृष्ठ स्तर समुच्चयों
$f(x, y, z) = c$ के रूप में भी उठ खड़े होते हैं; तब प्रवणक अभिलंब को दिशित करता है।

\section{प्राचलित पृष्ठ}

\begin{definition}[नियमित प्राचलित पृष्ठ]\label{def:b2:surfaces:param}
मान लीजिए $U \subseteq \R^2$ \omterm{def:b2:metric:topology}{विवृत} है। वर्ग $\mathcal{C}^k$ ($k \geq 1$) का \emph{प्राचलित
पृष्ठ}\index{प्राचलित पृष्ठ} वर्ग $\mathcal{C}^k$ का कोई प्रतिचित्रण $\sigma \colon U \to
\R^3$,
$(u, v) \mapsto \sigma(u, v)$, है। कोई बिंदु \emph{नियमित} है यदि आंशिक अवकलज सदिश
\[
\sigma_u(u, v) = \frac{\partial\sigma}{\partial u}(u,v),
\qquad
\sigma_v(u, v) = \frac{\partial\sigma}{\partial v}(u,v)
\]
रैखिकतः स्वतंत्र हों, अर्थात् $\sigma_u \wedge \sigma_v \neq 0$;
और पृष्ठ \emph{नियमित} है यदि उसका हर बिंदु नियमित हो।
\end{definition}

\begin{example}[तीन मानक वर्णन]\label{ex:b2:surfaces:standard}
\leavevmode
\begin{enumerate}
\item \emph{ग्राफ़:} $f \in
\mathcal{C}^1(U)$ के लिए $\sigma(u, v) = (u,\ v,\ f(u, v))$। सदा नियमित: $\sigma_u = (1, 0, f_u)$ तथा
$\sigma_v = (0, 1, f_v)$ स्वतंत्र हैं।
\item \emph{गोला (गोलीय निर्देशांक):} त्रिज्या $R$ के गोले के
लिए
\[
\sigma(\theta, \varphi)
= (R\cos\theta\cos\varphi,\ R\sin\theta\cos\varphi,\
R\sin\varphi),
\qquad
(\theta, \varphi) \in \R \times
\bigl(-\tfrac\pi2, \tfrac\pi2\bigr),
\]
जहाँ $\theta$ देशांतर है और $\varphi$ अक्षांश। जाँचने पर
$\norm{\sigma_\theta \wedge \sigma_\varphi} = R^2\cos\varphi > 0$:
अर्थात् ध्रुवों से दूर नियमित (जिन्हें यह चार्ट छोड़ देता है)।
\item \emph{स्तर समुच्चय:} $S = \{(x,y,z) : f(x, y, z) = c\}$, जहाँ
$f$ $\mathcal{C}^1$ है और $S$ पर $\nabla f \neq 0$। हर बिंदु के निकट
अंतर्निहित फलन प्रमेय (\cref{ch:b2:diffcalc}) से एक निर्देशांक शेष दो
के फलन के रूप में व्यक्त किया जा सकता है, अतः
$S$ स्थानीय रूप से कोई ग्राफ़ है।
\end{enumerate}
\end{example}

\begin{example}[स्तर समुच्चय से ग्राफ़ तक]
\label{ex:b2:surfaces:levelgraph}
मद 3 का अंतर्निहित फलन प्रमेय एक बार खुलकर चलाए जाने का हक़दार है।
गोला $x^2 + y^2 + z^2 = R^2$ उसके उत्तरी ध्रुव $(0, 0, R)$ के निकट
लीजिए: वहाँ $\frac{\partial f}{\partial z} = 2z =
2R \neq 0$, और $z$ के लिए हल करने पर ग्राफ़-चार्ट मिलता है
\[
z = \sqrt{R^2 - x^2 - y^2},
\qquad x^2 + y^2 < R^2 ,
\]
जो अपने (\omterm{def:b2:metric:topology}{विवृत}) प्रांत पर हर जगह नियमित है --- उस ध्रुव पर भी, जिसे
गोलीय चार्ट चूक गया था। $(R, 0, 0)$ जैसे किसी विषुवतीय बिंदु के
निकट वही प्रमेय इसके बदले $x$ के लिए हल कर देती है
($\frac{\partial f}{\partial x} = 2R \neq 0$)। अंगूठे का नियम: कोई स्तर पृष्ठ उस
निर्देशांक समतल के ऊपर ग्राफ़ होता है जो \emph{प्रवणक के सबसे बड़े
घटक के लांबिक} है; और गोले को ऐसे छह ग्राफ़-चार्टों से ढँककर उसकी
चिकनाई हर जगह बिना ज़रा भी त्रिकोणमिति के जाँची जा सकती है।
\end{example}

\section{स्पर्श समतल और अभिलंब}

\begin{definition}[स्पर्श समतल]\label{def:b2:surfaces:tangent}
मान लीजिए $\sigma$ $(u_0, v_0)$ पर नियमित है, $M_0 = \sigma(u_0, v_0)$।
\emph{स्पर्श समतल}\index{स्पर्श समतल} $T_{M_0}S$ वह समतल है जो $M_0$ से होकर
जाता है और $\operatorname{Vect}(\sigma_u, \sigma_v)$ से दिशित है ($(u_0, v_0)$ पर आंशिक
अवकलज)। \emph{इकाई अभिलंब}\index{इकाई अभिलंब} है
\[
n(u_0, v_0)
= \frac{\sigma_u \wedge \sigma_v}
       {\norm{\sigma_u \wedge \sigma_v}} .
\]
\end{definition}

\begin{proposition}[स्पर्शी सदिश वेग सदिश ही
हैं]\label{prop:b2:surfaces:velocity}
$T_{M_0}S$ की दिशा ठीक सदिशों $\gamma'(0)$ का समुच्चय है, जहाँ
$\gamma = \sigma \circ c$ $M_0$ से होकर पृष्ठ पर खींचे गए $\mathcal{C}^1$ वक्रों पर
घूमता है (अर्थात् $c \colon (-\varepsilon, \varepsilon) \to U$ $\mathcal{C}^1$ है, जहाँ $c(0) = (u_0, v_0)$)।
\end{proposition}

\begin{proof}
यदि $c(t) = (u(t), v(t))$, तो शृंखला नियम
(\cref{ch:b2:diffcalc}) देता है
\[
\gamma'(0) = u'(0)\,\sigma_u + v'(0)\,\sigma_v
\in \operatorname{Vect}(\sigma_u, \sigma_v).
\]
विलोमतः, सदिश $a\sigma_u + b\sigma_v$ वक्र $c(t) = (u_0 + at,\ v_0 + bt)$ से प्राप्त होता है, जो
छोटे $\abs t$ के लिए \omterm{def:b2:metric:topology}{विवृत} समुच्चय $U$ में बना रहता है।
\end{proof}

\begin{example}[हेलिकॉइड का स्पर्श समतल]
\label{ex:b2:surfaces:helicoidtangent}
हेलिकॉइड $\sigma(u, v) = (v\cos u,\ v\sin u,\ au)$ के लिए,
बिंदु $\sigma(0, 1) = (1, 0, 0)$ पर:
\[
\sigma_u = (0,\ 1,\ a), \qquad \sigma_v = (1,\ 0,\ 0),
\qquad
\sigma_u \wedge \sigma_v = (0,\ a,\ -1),
\]
अतः \omterm{def:b2:surfaces:tangent}{स्पर्श समतल} $a\,y = z$ है। उसमें पूरा क्षैतिज रेखक $\{(t, 0, 0)\}$
समाया है (दिशा $\sigma_v$): \cref{exo:b2:surfaces:1} के शंकु की तरह, सरल
रेखाओं से रेखित कोई भी पृष्ठ अपने हर रेखक को उसके अनुदिश \omterm{def:b2:surfaces:tangent}{स्पर्श
समतल} के भीतर रखता है। दूसरी स्पर्शी दिशा $\sigma_u$ कुंडली $u \mapsto \sigma(u, 1)$
का वेग है: एक चार्ट, दो खींचे गए वक्र, और पूरा \omterm{def:b2:surfaces:tangent}{स्पर्श समतल} \omterm{def:b2:structures:generated}{जनित} ---
\cref{prop:b2:surfaces:velocity} काम करती हुई।
\end{example}

\begin{proposition}[स्तर पृष्ठ का अभिलंब]\label{prop:b2:surfaces:gradient}
मान लीजिए $S = \{f = c\}$, जहाँ $f$ वर्ग $\mathcal{C}^1$ का है और $\nabla
f(M_0) \neq 0$। तब $M_0$ पर
$S$ का \omterm{def:b2:surfaces:tangent}{स्पर्श समतल} वह समतल है जो $M_0$ से होकर जाता है और
$\nabla f(M_0)$ के लांबिक है:
\[
T_{M_0}S :\quad
\langle \nabla f(M_0),\ M - M_0\rangle = 0 .
\]
\end{proposition}

\begin{proof}
$M_0$ से होकर $S$ पर खींचे गए किसी भी वक्र $\gamma$ के लिए
सर्वसमिक रूप से $f(\gamma(t)) =
c$, अतः शृंखला नियम $\langle
\nabla f(M_0), \gamma'(0)\rangle = 0$ देता है: यानी सारे
वेग सदिश प्रवणक के लांबिक हैं, इसलिए स्पर्शी दिशा समतल $\nabla f(M_0)^\perp$ में
समाई है। दोनों $2$-विमीय उपसमष्टियाँ हैं --- स्पर्शी दिशा इसलिए कि
$S$ स्थानीय रूप से कोई नियमित ग्राफ़ है
(\cref{ex:b2:surfaces:standard}), और लांबिक पूरक इसलिए कि $\nabla f(M_0) \neq 0$ ---
अतः वे बराबर हैं।
\end{proof}

\begin{example}\label{ex:b2:surfaces:spheretangent}
गोले $x^2 + y^2 + z^2 = R^2$ के लिए: $\nabla f = 2(x, y, z)$,
अतः $M_0$ पर \omterm{def:b2:surfaces:tangent}{स्पर्श समतल} त्रिज्या $\vect{OM_0}$ के लांबिक है --- यही
वह शास्त्रीय तथ्य है कि त्रिज्या और \omterm{def:b2:surfaces:tangent}{स्पर्श समतल} परस्पर लंब होते हैं,
और समीकरण $\langle M_0, M\rangle = R^2$ है।
\end{example}

\begin{example}[ग्राफ़ का स्पर्श समतल]\label{ex:b2:surfaces:graphtangent}
$(x_0, y_0)$ पर $z = f(x, y)$ के लिए: $F(x,y,z) = f(x,y) - z$ पर
\cref{prop:b2:surfaces:gradient} लगाने पर,
\[
z = f(x_0, y_0)
+ f_x(x_0, y_0)(x - x_0)
+ f_y(x_0, y_0)(y - y_0) ,
\]
अर्थात् प्रथम-कोटि टेलर प्रसार का \omterm{def:b2:affine:subspace}{आनत} भाग --- \omterm{def:b2:surfaces:tangent}{स्पर्श समतल} \omterm{def:b2:diffcalc:differential}{अवकल} का
ग्राफ़ है, जैसा होना ही चाहिए।
\end{example}

\begin{example}[पृष्ठ का निकटतम
बिंदु]\label{ex:b2:surfaces:closest}
परवलयज $z = x^2 + y^2$ का कौन-सा बिंदु $P = (0, 0, 1)$ के सबसे निकट है?
पृष्ठ के अनुदिश दूरी के वर्ग को न्यूनतम कीजिए: $\rho^2 = x^2 + y^2$ के साथ,
\[
g(\rho^2) = \rho^2 + (\rho^2 - 1)^2,
\qquad
g'(\rho^2) = 1 + 2(\rho^2 - 1) = 0
\iff \rho^2 = \tfrac12 ,
\]
जिससे ऊँचाई $z = \frac12$ तथा दूरी $\sqrt{\tfrac12 + \tfrac14} = \frac{\sqrt3}2$ वाले
बिंदुओं का वृत्त मिलता है। न्यूनतमता का ज्यामितीय हस्ताक्षर: ऐसे किसी
बिंदु $M$ पर सदिश $\vect{MP}$ को पृष्ठ का \emph{अभिलंब} होना ही
चाहिए --- अन्यथा किसी खींचे गए वक्र के अनुदिश ऐसे वेग से सरकना जिसका
कोई घटक $P$ की ओर हो, दूरी घटा देता। जाँच: $M = (x, y,
\tfrac12)$ पर
$\nabla(z - x^2 - y^2) = (-2x, -2y, 1)$, जबकि $\vect{MP} = (-x, -y, \tfrac12) =
\tfrac12(-2x, -2y, 1)$: समांतर, जैसा अनुमानित था। ``लंब का
पाद'' वाली यही प्रथम-कोटि शर्त \cref{exo:b2:surfaces:12} में स्तर समुच्चयों पर
चरम मानों को भी चलाएगी।
\end{example}

\begin{example}[द्विघातजों के स्पर्श समतल: ध्रुवण
नियम]\label{ex:b2:surfaces:polarization}
मान लीजिए $S : xy + yz + zx = 1$ और $M_0 = (1, 1, 0) \in S$। यहाँ
$\nabla f = (y + z,\ x + z,\ x + y)$, अतः $\nabla f(M_0) = (1,
1, 2)$ और \omterm{def:b2:surfaces:tangent}{स्पर्श समतल} है
\[
(x - 1) + (y - 1) + 2z = 0,
\qquad\text{अर्थात्}\qquad x + y + 2z = 2 .
\]
वही उत्तर \emph{\omterm{def:b2:quadratic:def}{ध्रुवण} नियम} से भी आता है, जो \cref{ex:b2:surfaces:spheretangent} तथा
\cref{exo:b2:surfaces:2} का व्यापकीकरण है: \omterm{pb:b2:surfaces:1}{द्विघातज} के समीकरण
में $x^2$ के स्थान पर $x_0x$ रखिए, और हर गुणनफल $xy$ के
स्थान पर $\frac{x_0y +
y_0x}2$ (और \omterm{def:b2:structures:generated}{चक्रीय} रूप से आगे भी):
\[
\frac{x_0y + y_0x}2 + \frac{y_0z + z_0y}2 + \frac{z_0x +
x_0z}2 = 1
\;\xrightarrow{\ M_0 = (1,1,0)\ }\;
\frac{x + y}2 + \frac z2 + \frac z2 = 1 ,
\]
जो फिर वही $x + y + 2z = 2$ है। नियम इसलिए चलता है क्योंकि किसी \omterm{def:b2:quadratic:def}{द्विघात रूप}
का $\nabla$ आधार बिंदु के सामने मूल्यांकित \omterm{def:b2:metric:connected}{संबद्ध} द्विरैखिक रूप ही है
--- किसी \omterm{pb:b2:surfaces:1}{द्विघातज} को स्पर्श करना \omterm{def:b2:quadratic:def}{ध्रुवण} है, \cref{ch:b2:quadratic} का एक और चेहरा।
\end{example}

\section{प्रथम आधारभूत रूप}

\begin{definition}[प्रथम आधारभूत रूप]\label{def:b2:surfaces:fff}
मान लीजिए $\sigma \colon U \to \R^3$ कोई नियमित $\mathcal{C}^1$ पृष्ठ है। $(u,v)$ पर उसका
\emph{प्रथम आधारभूत रूप}\index{प्रथम आधारभूत रूप} $\R^2$ पर धनात्मक
निश्चित \omterm{def:b2:quadratic:def}{द्विघात रूप} है
\[
I(h, k) = \norm{h\,\sigma_u + k\,\sigma_v}^2
= E\,h^2 + 2F\,hk + G\,k^2,
\]
जहाँ
\[
E = \norm{\sigma_u}^2,
\qquad
F = \langle \sigma_u, \sigma_v\rangle,
\qquad
G = \norm{\sigma_v}^2 .
\]
\end{definition}

\begin{remark}
$I$ परिवेशी यूक्लिडीय अंतर्गुणन का \omterm{def:b2:surfaces:tangent}{स्पर्श समतल} तक प्रतिबंधन है,
जिसे आधार $(\sigma_u, \sigma_v)$ में पढ़ा गया है: वह ठीक इसलिए धनात्मक
निश्चित है क्योंकि $\sigma_u, \sigma_v$ स्वतंत्र हैं
(\cref{ch:b2:quadratic})। पृष्ठ पर हर दूरिक राशि --- खींचे गए वक्रों
की लंबाइयाँ, उनके बीच के कोण, \omterm{def:b2:surfaces:area}{क्षेत्रफल} --- अकेले $E, F, G$ से परिकलित
होती है। उपयुक्त प्राचलों में एक ही $E,
F, G$ वाले दो पृष्ठ
\emph{समदूरिक} हैं, चाहे वे अवकाश में अलग-अलग बैठे हों: आंतरिक
ज्यामिति का आरंभ बिंदु यही है।
\end{remark}

\begin{example}[निर्देशांक वक्रों के बीच के
कोण]\label{ex:b2:surfaces:coordangle}
\omterm{def:b2:surfaces:fff}{प्रथम आधारभूत रूप} कोण भी मापता है: निर्देशांक वक्र $u \mapsto \sigma(u, v_0)$ और $v \mapsto
\sigma(u_0, v)$
कोण $\theta$ पर मिलते हैं, जहाँ
\[
\cos\theta = \frac{\langle\sigma_u,
\sigma_v\rangle}{\norm{\sigma_u}\,\norm{\sigma_v}}
= \frac{F}{\sqrt{EG}} :
\]
अर्थात् अकेला गुणांक $F$ प्राचल-जाल की लांबिकता तय कर देता है।
गोलीय चार्ट तथा हेलिकॉइड के लिए $F =
0$: याम्योत्तर अक्षांश-वृत्तों
को, और कुंडलियाँ क्षैतिज रेखकों को, समकोण पर काटती हैं --- और इसीलिए
उनके क्षेत्रफल-समाकल्य $\sqrt{EG}$ तक सिमट गए थे। किसी ग्राफ़-चार्ट के
लिए $F = f_xf_y$ केवल वहीं लुप्त होता है जहाँ कोई आंशिक अवकलज लुप्त हो:
किसी झुके हुए ग्राफ़ का निर्देशांक-जाल लांबिक \emph{नहीं} होता, यद्यपि
नीचे वाला $(x, y)$-जाल होता है। जब किसी पृष्ठ पर परिकलन भारी दिखने
लगें, तब पहली चाल $F = 0$ वाला चार्ट खोजना है।
\end{example}

\begin{example}[काठी वाला चार्ट]\label{ex:b2:surfaces:saddlechart}
काठी $z = xy$ के लिए चार्ट $\sigma(u, v) = (u, v,
uv)$ के साथ:
\[
\sigma_u = (1, 0, v), \qquad \sigma_v = (0, 1, u),
\qquad
E = 1 + v^2, \quad F = uv, \quad G = 1 + u^2 ,
\]
और $EG - F^2 = 1 + u^2 + v^2 > 0$: अर्थात् हर जगह नियमित। किसी बिंदु से होकर जाने वाले
दोनों निर्देशांक वक्र $\R^3$ की \emph{सरल रेखाएँ} हैं ($u$ स्थिर
कीजिए या $v$ स्थिर कीजिए: द्विगुण रेखित काठी के रेखक), फिर भी
अक्षों से हटकर $F \neq 0$: किसी व्यापक बिंदु से होकर जाने वाले रेखक
लांबिक नहीं होते। दोनों रेखक \omterm{def:b2:surfaces:tangent}{स्पर्श समतल} में पड़ते हैं और उसे \omterm{def:b2:structures:generated}{जनित}
करते हैं --- अतः \omterm{def:b2:surfaces:tangent}{स्पर्श समतल} पृष्ठ को दो पूरी रेखाओं के अनुदिश काटता
है, जो गोले का ठीक उल्टा छोर है, क्योंकि गोले के \omterm{def:b2:surfaces:tangent}{स्पर्श समतल} केवल एक
बिंदु पर छूते हैं। किसी पृष्ठ और उसके स्पर्श समतलों के बीच के
``द्वितीय-कोटि संपर्क'' का चिह्न \omterm{thm:b2:curves:frenet2d}{वक्रता} की कहानी है, जिसे तृतीय वर्ष
का खंड उठाता है।
\end{example}

\begin{proposition}[पृष्ठ पर खींचे गए वक्र की
लंबाई]\label{prop:b2:surfaces:curvelength}
यदि $\gamma(t) = \sigma(u(t), v(t))$, $t \in [a, b]$, $\mathcal{C}^1$ हो, तो
\[
L(\gamma) = \int_a^b
\sqrt{E\,u'^2 + 2F\,u'v' + G\,v'^2}\;\dd t ,
\]
जहाँ $E, F, G$ का मूल्यांकन $(u(t), v(t))$ पर होता है।
\end{proposition}

\begin{proof}
शृंखला नियम से $\gamma' = u'\sigma_u + v'\sigma_v$, अतः
$\norm{\gamma'}^2 = I(u', v') = Eu'^2 + 2Fu'v' + Gv'^2$; $\norm{\gamma'}$ का समाकलन कीजिए
(\cref{def:b2:curves:length})।
\end{proof}

\begin{example}[हवाई जहाज़ ध्रुव के ऊपर से क्यों उड़ते हैं]\label{ex:b2:surfaces:polarroute}
दो हवाई अड्डे अक्षांश $\varphi_0$ पर और विपरीत देशांतरों पर बैठे हैं:
त्रिज्या $R$ के गोले पर $A = \sigma(0, \varphi_0)$ तथा $B = \sigma(\pi,
\varphi_0)$। अक्षांश-वृत्त के
अनुदिश ($\varphi \equiv \varphi_0$) \omterm{def:b2:curves:length}{लंबाई}
$\int_0^\pi R\cos\varphi_0\,\dd\theta = \pi R\cos\varphi_0$ है। ध्रुव के ऊपर वाले मार्ग के अनुदिश (याम्योत्तर $\theta = 0$ से
ऊपर, याम्योत्तर $\theta = \pi$ से नीचे) वह $2R(\frac\pi2 -
\varphi_0)$ है। अक्षांश $\varphi_0 = \frac\pi3$ (साठ
अंश) पर: अक्षांश-मार्ग $\pi R/2 \approx 1.571\,R$, ध्रुवीय मार्ग
$\pi R/3 \approx 1.047\,R$ --- एक तिहाई छोटा। वस्तुतः $\intcc0{\pi/2}$ पर
$\pi\cos\varphi_0 \geq \pi - 2\varphi_0$ (फलन $\pi\cos\varphi - \pi +
2\varphi$ दोनों सिरों पर लुप्त होता है और उसका अवकलज
$2 -
\pi\sin\varphi$ एक बार चिह्न बदलता है, अतः वह पहले वर्धमान फिर ह्रासमान है,
इसलिए अऋणात्मक): यानी ध्रुवीय मार्ग कभी नहीं हारता। \omterm{def:b2:surfaces:fff}{प्रथम आधारभूत
रूप} ने किसी नौवहन प्रश्न को दो एक-पंक्ति समाकलों में बदल दिया;
\cref{exo:b2:surfaces:6} इस विचार को याम्योत्तरों के लिए एक सच्ची न्यूनतमता-उपपत्ति
तक ले जाता है।
\end{example}

\begin{lemma}[लाग्रांज सर्वसमिका]\label{lem:b2:surfaces:lagrange}
सब $a, b \in \R^3$ के लिए:
$\norm{a \wedge b}^2 = \norm a^2 \norm b^2 - \langle a, b\rangle^2$।
विशेष रूप से
\[
\norm{\sigma_u \wedge \sigma_v} = \sqrt{EG - F^2}.
\]
\end{lemma}

\begin{proof}
यदि हम $b$ के स्थान पर $a$ के लांबिक उसका घटक $b_\perp
= b - \frac{\langle a, b\rangle}{\norm a^2}a$ रख दें
तो दोनों पक्ष अपरिवर्तित रहते हैं ($a \neq 0$ के लिए; स्थिति $a = 0$ तुच्छ
है): बायाँ पक्ष इसलिए कि $a \wedge a = 0$, और दायाँ पक्ष $\norm
{b_\perp}^2 = \norm b^2 - \frac{\langle a,b\rangle^2}{\norm a^2}$ तथा
$\langle a, b_\perp\rangle = 0$ को खोलकर। अतः सर्वसमिका को केवल लांबिक $a, b$ के लिए सिद्ध
करना पर्याप्त है, जहाँ वह $\norm{a \wedge
b} = \norm a \norm b$ पढ़ी जाती है: और यह सत्य है,
क्योंकि लांबिक सदिशों के लिए सदिश गुणनफल का \omterm{def:b2:nvs:norm}{मानदंड} $\norm a\norm b\,\abs{\sin\frac\pi2}$ होता है।
प्रदर्शित सूत्र स्थिति $a = \sigma_u$, $b = \sigma_v$ है।
\end{proof}

\begin{remark}
लाग्रांज सर्वसमिका कहती है कि $EG - F^2 = \det\begin{psmallmatrix}
E & F\\ F & G\end{psmallmatrix}$ $(\sigma_u, \sigma_v)$ का \emph{ग्राम \omterm{def:b2:linalg:det}{सारणिक}}
है: अर्थात् उनसे \omterm{def:b2:structures:generated}{जनित} समांतर चतुर्भुज के \omterm{def:b2:surfaces:area}{क्षेत्रफल} का वर्ग।
नियमितता, \omterm{def:b2:surfaces:fff}{प्रथम आधारभूत रूप} की धनात्मक निश्चितता, और ग्राम \omterm{def:b2:linalg:det}{सारणिक} की
धनात्मकता एक ही शर्त के तीन कथन हैं --- और इसीलिए नीचे दिया
क्षेत्रफल-समाकल्य किसी नियमित चार्ट पर कभी लुप्त नहीं होता।
\end{remark}

\begin{definition}[क्षेत्रफल]\label{def:b2:surfaces:area}
मान लीजिए $\sigma \colon U \to \R^3$ कोई नियमित एकैकी $\mathcal{C}^1$ पृष्ठ है और $K \subseteq U$ कोई \omterm{def:b2:metric:compact}{संहत}
प्रांत, जिस पर द्विक समाकल अर्थपूर्ण हों (\cref{ch:b2:multint})। टुकड़े $\sigma(K)$ का
\emph{क्षेत्रफल}\index{क्षेत्रफल!पृष्ठ का} है
\[
\mathcal{A}
= \iint_K \norm{\sigma_u \wedge \sigma_v}\,\dd u\,\dd v
= \iint_K \sqrt{EG - F^2}\;\dd u\,\dd v .
\]
\end{definition}

\begin{remark}[यह सूत्र क्यों]
आयत $[u, u + \dd u] \times [v, v + \dd v]$ प्रथम कोटि तक $\sigma_u\,\dd u$ और $\sigma_v\,\dd v$ से \omterm{def:b2:structures:generated}{जनित} समांतर
चतुर्भुज पर भेजा जाता है, जिसका \omterm{def:b2:surfaces:area}{क्षेत्रफल} $\norm{\sigma_u \wedge
\sigma_v}\,\dd u\,\dd v$ है: परिभाषा
स्थानीय क्षेत्रफल-विरूपण गुणक का समाकलन करती है, ठीक वैसे जैसे
\omterm{def:b2:curves:length}{चाप-लंबाई} स्थानीय चाल का समाकलन करती है। प्राचल-परिवर्तन के साथ
संगति \cref{exo:b2:surfaces:8} है; द्विक समाकलों के चर-परिवर्तन सूत्र के साथ संगति
\cref{ch:b2:multint} में चर्चित है।
\end{remark}

\begin{example}[दो अलग ग्राफ़, एक ही
क्षेत्रफल]\label{ex:b2:surfaces:samearea}
इकाई चक्रिका के ऊपर कटोरे $z = \frac12(x^2 + y^2)$ और काठी
$z = xy$ की तुलना कीजिए। उनके क्षेत्रफल-समाकल्य
(\cref{exo:b2:surfaces:5}) हैं
\[
\sqrt{1 + x^2 + y^2}
\qquad\text{तथा}\qquad
\sqrt{1 + y^2 + x^2} :
\]
अर्थात् एक ही। दोनों पृष्ठ --- एक हर दिशा में एक ही ओर मुड़ता हुआ,
दूसरा काठी के आकार का --- हर प्रांत पर ठीक बराबर \omterm{def:b2:surfaces:area}{क्षेत्रफल} रखते हैं,
इकाई चक्रिका के ऊपर $\frac{2\pi}3(2\sqrt2 - 1)$। क्षेत्रफल-अवयव केवल प्रवणक की
\emph{\omterm{def:b2:curves:length}{लंबाई}} देखता है, मुड़ाव की व्यवस्था नहीं; कटोरे को काठी से
अलग पहचानने के लिए द्वितीय-कोटि आँकड़े चाहिए (\cref{fig:b2:surfaces:saddle} में दिखाई
चिह्न-संरचना), जिन्हें कितना भी क्षेत्रफल-मापन नहीं पकड़ पाता। \omterm{def:b2:surfaces:fff}{प्रथम
आधारभूत रूप}: दूरिक, आकार के प्रति अंधा; आकार देखने वाला द्वितीय रूप
तृतीय वर्ष का है।
\end{example}

\begin{example}[गोले का क्षेत्रफल]\label{ex:b2:surfaces:spherearea}
\cref{ex:b2:surfaces:standard} के गोलीय चार्ट के लिए:
\[
\sigma_\theta = R(-\sin\theta\cos\varphi,\
\cos\theta\cos\varphi,\ 0),
\qquad
\sigma_\varphi = R(-\cos\theta\sin\varphi,\
-\sin\theta\sin\varphi,\ \cos\varphi),
\]
अतः $E = R^2\cos^2\varphi$, $F = 0$, $G = R^2$ और $\sqrt{EG - F^2}
= R^2\cos\varphi$। इसलिए
\[
\mathcal{A}
= \int_{-\pi/2}^{\pi/2}\!\!\int_0^{2\pi}
R^2\cos\varphi\;\dd\theta\,\dd\varphi
= 2\pi R^2\,\bigl[\sin\varphi\bigr]_{-\pi/2}^{\pi/2}
= \boxed{4\pi R^2} .
\]
\end{example}

\begin{example}[शंकु, स्कूली सूत्र के सामने
जाँचा गया]\label{ex:b2:surfaces:conearea}
वलय $a \leq
\rho \leq b$ के ऊपर शंकु $z = \sqrt{x^2 + y^2}$ के लिए \cref{exo:b2:surfaces:5} का ग्राफ़-सूत्र
$1 + f_x^2 + f_y^2 = 2$ देता है (परिकलित कीजिए $f_x = x/\rho$, $f_y = y/\rho$), अतः
\[
\mathcal A = \sqrt2\,\pi\,(b^2 - a^2) .
\]
तिर्यक-ऊँचाई सूत्र $\pi\rho\ell$ (\cref{ex:b2:surfaces:revolution} का) के सामने संगति-जाँच: आधार
त्रिज्या $b$ तथा $a$ वाले पूरे शंकुओं के पार्श्व \omterm{def:b2:surfaces:area}{क्षेत्रफल}
$\pi b\cdot b\sqrt2$ तथा $\pi a\cdot a\sqrt2$ हैं, जिनका अंतर ठीक $\sqrt2\pi(b^2 - a^2)$ है। दो चार्ट, दो
सूत्र, एक ही \omterm{def:b2:surfaces:area}{क्षेत्रफल} --- \cref{exo:b2:surfaces:8} में सिद्ध अपरिवर्त्यता, खुले में
देखी हुई।
\end{example}

\begin{remark}[क्षेत्रफलों के लिए त्वरित जाँचें]
तीन तात्क्षणिक जाँचें क्षेत्रफल-परिकलन की अधिकांश भूलें पकड़ लेती
हैं। \emph{मापन}: किसी पृष्ठ को $\lambda$ से फैलाने पर $E, F, G$ $\lambda^2$ से
गुणित हो जाता है और \omterm{def:b2:surfaces:area}{क्षेत्रफल} $\lambda^2$ से --- जिस उत्तर की $R$ पर
निर्भरता वर्गिक न हो (जैसे $4\pi R^2$), वह ग़लत है।
\emph{अवयव की धनात्मकता}: चार्ट के अभ्यंतर पर $\sqrt{EG - F^2}$ दृढ़ता से
धनात्मक होना ही चाहिए; कोई लुप्त होता मान चार्ट के अपभ्रंश का संकेत
है, जिसे गोले के ध्रुवों की तरह काट देना चाहिए। \emph{सममिति}: किसी
\omterm{def:b2:quadratic:adjoint}{सममित} टुकड़े पर परिकलन उसके सर्वांगसम भागों को जोड़ने के साथ संगत
होना चाहिए --- अर्धगोले को $2\pi R^2$ देना ही चाहिए।
\end{remark}

\begin{example}[घूर्णन पृष्ठ]\label{ex:b2:surfaces:revolution}
वर्ग $\mathcal{C}^1$ के वक्र $z \mapsto (r(z), 0, z)$, $r > 0$, को $z$-अक्ष के परितः घुमाइए:
\[
\sigma(\theta, z) = (r(z)\cos\theta,\ r(z)\sin\theta,\ z).
\]
तब $E = r(z)^2$, $F = 0$, $G = 1 + r'(z)^2$, अतः
\[
\mathcal{A} = \int_a^b\!\!\int_0^{2\pi}
r(z)\sqrt{1 + r'(z)^2}\;\dd\theta\,\dd z
= 2\pi\int_a^b r(z)\sqrt{1 + r'(z)^2}\,\dd z ,
\]
यही शास्त्रीय सूत्र है (परिधि $2\pi r$ गुणा तिर्यक लंबाई-अवयव)। शंकु
$r(z) = kz$, $z \in [0, h]$, के लिए: $\mathcal
A = 2\pi k\sqrt{1 + k^2}\,\frac{h^2}{2} = \pi \rho \ell$, जहाँ $\rho = kh$ आधार त्रिज्या है और
$\ell = h\sqrt{1 + k^2}$ तिर्यक ऊँचाई --- वही स्कूली सूत्र, जो अब मान लिया नहीं, बल्कि
व्युत्पन्न किया गया है।
\end{example}

\begin{example}[कैटेनॉइड]\label{ex:b2:surfaces:catenoid}
कैटेनरी $r(z) = \cosh z$, $z \in \intcc{-1}{1}$, को उसकी अक्ष के परितः घुमाइए: परिणामी
\emph{कैटेनॉइड} का, घूर्णन सूत्र तथा $1 + \sinh^2 z = \cosh^2 z$ से,
\[
\mathcal A = 2\pi\int_{-1}^1\cosh z\,\sqrt{1 + \sinh^2z}\,
\dd z = 2\pi\int_{-1}^1\cosh^2z\,\dd z
= \pi\bigl[z + \sinh z\cosh z\bigr]_{-1}^{1},
\]
अर्थात्
\[
\mathcal A = 2\pi + \pi\sinh 2 \approx 17.68 .
\]
तिर्यक गुणक $\sqrt{1 + r'^2}$ त्रिज्या के साथ मिलकर पूर्ण वर्ग बन गया --- वही
सर्वसमिका जिसने वक्र वाले अध्याय में कैटेनरी की \omterm{def:b2:curves:length}{चाप-लंबाई} को
प्रारंभिक बना दिया था। यह बीजगणित का संयोग नहीं है: दोनों सीमांत
वृत्तों को घेरने वाले सारे घूर्णन पृष्ठों में कैटेनॉइड \omterm{def:b2:surfaces:area}{क्षेत्रफल}
न्यूनतम करता है (वह दो छल्लों के बीच की साबुन-झिल्ली का आकार है), और
यही विचरणीय गुणधर्म $\cosh$ को अलग से चुन लेता है; तृतीय वर्ष का खंड
इसे विचरण-कलन से सिद्ध करता है।
\end{example}

\begin{figure}[ht]
\centering
\begin{tikzpicture}[scale=1.05]
  % A saddle surface z = x^2 - y^2 drawn as a wire mesh in oblique projection
  % projection: (x,y,z) -> (x + 0.45*y, z + 0.35*y)
  \begin{scope}
    % u-lines (x varies, y fixed)
    \foreach \yy in {-1,-0.5,0,0.5,1} {
      \draw[ocDarkBlue!70, thick, domain=-1:1, smooth, samples=25]
        plot ({\x + 0.45*\yy}, {\x*\x - \yy*\yy + 0.35*\yy});
    }
    % v-lines (y varies, x fixed)
    \foreach \xx in {-1,-0.5,0,0.5,1} {
      \draw[ocMint!80!black, thick, domain=-1:1, smooth, samples=25]
        plot ({\xx + 0.45*\x}, {\xx*\xx - \x*\x + 0.35*\x});
    }
    % tangent plane at origin: z = 0 -> parallelogram
    \draw[ocRed, dashed]
      (-0.8 - 0.36, -0.28) -- (0.8 - 0.36, -0.28)
      -- (0.8 + 0.36, 0.28) -- (-0.8 + 0.36, 0.28) -- cycle;
    \fill (0,0) circle (0.045) node[below right, font=\small] {$M_0$};
    % normal vector at origin
    \draw[->, thick] (0,0) -- (0,0.9)
      node[right, font=\small] {$n$};
  \end{scope}
\end{tikzpicture}
\caption{मूल बिंदु के निकट काठी $z = x^2 - y^2$, उसके निर्देशांक वक्रों
($u$-वक्र नीले, $v$-वक्र हरे), $M_0 = (0,0,0)$ पर \omterm{def:b2:surfaces:tangent}{स्पर्श समतल} (बिंदुकित)
तथा \omterm{def:b2:surfaces:tangent}{इकाई अभिलंब} $n$ के साथ। पृष्ठ अपने \omterm{def:b2:surfaces:tangent}{स्पर्श समतल} को पार करता
है --- नति-परिवर्तन का द्विविमीय प्रतिरूप।}
\label{fig:b2:surfaces:saddle}
\end{figure}

\begin{remark}[सामान्य चूकें]
(क) \emph{चार्ट की विचित्रताएँ पृष्ठ की विचित्रताएँ नहीं होतीं}:
गोलीय चार्ट ध्रुवों पर अपभ्रष्ट हो जाता है ($\cos\varphi = 0$), पर गोला वहाँ
पूरी तरह चिकना है --- कोई दूसरा चार्ट (अक्षों की भूमिकाएँ बदल दीजिए)
ध्रुवों पर नियमित है। किसी बिंदु को विचित्र घोषित करने से पहले दूसरा
प्राचलन आज़माइए। (ख) \emph{$\sigma$ की नियमितता प्राचलन से संबंध रखती
है, प्रतिबिंब से नहीं}: $\sigma(u, v) = (u^3, v,
0)$ $u = 0$ के अनुदिश अनियमित है, यद्यपि
उसका प्रतिबिंब कोई समतल है। (ग) क्षेत्रफल-सूत्र माँगता है कि $\sigma$
$K$ पर \emph{एकैकी} हो: किसी टुकड़े को दो बार ढँकने वाला चार्ट
उसे दो बार गिन लेता है ($\theta$ का $\intcc0{4\pi}$ पर चलना गोले का \omterm{def:b2:surfaces:area}{क्षेत्रफल}
दुगुना कर देता है)। (घ) \omterm{def:b2:surfaces:tangent}{इकाई अभिलंब} पृष्ठ द्वारा चिह्न तक परिभाषित
होता है, पर उसे \emph{चुनता} चार्ट है ($u, v$ का क्रम); अभिविन्यास
वाले कथनों को वह चुनाव नियत करना ही चाहिए। (ङ) अंत में, $EG - F^2 > 0$ कोई
अतिरिक्त परिकल्पना नहीं है: लाग्रांज सर्वसमिका से वह ठीक नियमितता
ही है --- यदि वह कहीं लुप्त हो जाए, तो दोष चार्ट का है, और वहाँ न
कोई \omterm{def:b2:surfaces:area}{क्षेत्रफल} सूत्र लागू होता है, न कोई स्पर्श-समतल सूत्र।
\end{remark}

\begin{remark}[इस खंड के भीतर के परिप्रेक्ष्य]
यहाँ से आगे की कड़ियाँ। क्षेत्रफल-अवयव $\sqrt{EG -
F^2}\,\dd u\,\dd v$ भेस बदले हुए
द्विविमीय याकोबी है, और \cref{ch:b2:multint} उस उपमा को चर-परिवर्तन प्रमेय के
साथ ठीक-ठीक बना देता है --- वहाँ के पृष्ठ समाकल इसी अध्याय के
\omterm{def:b2:surfaces:area}{क्षेत्रफल} हैं, बस ऊपर से एक समाकल्य और। \omterm{def:b2:surfaces:fff}{प्रथम आधारभूत रूप} धनात्मक
द्विघात रूपों का कोई क्षेत्र है, जिसे \cref{ch:b2:quadratic} के औज़ार \omterm{def:b2:funcseq:def}{बिंदुवार} साध
लेते हैं, और जिसकी वर्णक्रमीय प्रमेय इस अध्याय के सप्ताहांत वाले
द्विघातज-वर्गीकरण को भी चलाती है। और अभिलंब रेखा प्रतिबंध-समुच्चयों
पर चरम-मान समस्याओं को चलाती है (\cref{ex:b2:surfaces:closest}), जो \cref{thm:b2:diffcalc:extrema} के साथ रेखांकित
लाग्रांज गुणक विधि का ज्यामितीय बीज है।
\end{remark}

\begin{omfigure}
\begin{tikzpicture}[scale=0.8]
  \draw[omDef, thick] (0,0) ellipse (1.5 and 0.95);
  \draw[omThm] (0,0) ellipse (1.5 and 0.35);
  \node[below, font=\small] at (0,-1.25) {दीर्घवृत्तज};
\end{tikzpicture}
\qquad
\begin{tikzpicture}[scale=0.8]
  \draw[omDef, thick] plot[domain=-1.1:1.1, samples=30]
    ({0.6*(exp(\x)+exp(-\x))/2}, {\x});
  \draw[omDef, thick] plot[domain=-1.1:1.1, samples=30]
    ({-0.6*(exp(\x)+exp(-\x))/2}, {\x});
  \draw[omThm] (0,0) ellipse (0.6 and 0.16);
  \draw[omDef] (0,1.1) ellipse (1.0 and 0.22);
  \draw[omDef] (0,-1.1) ellipse (1.0 and 0.22);
  \node[below, font=\small] at (0,-1.55)
    {एकपृष्ठी अतिपरवलयज};
\end{tikzpicture}
\qquad
\begin{tikzpicture}[scale=0.8]
  \draw[omDef, thick] plot[domain=-1.05:1.05, samples=30]
    ({\x}, {1.1*\x*\x - 1.1});
  \draw[omThm] (0,0.055) ellipse (1.0 and 0.22);
  \node[below, font=\small] at (0,-1.45)
    {दीर्घवृत्तीय परवलयज};
\end{tikzpicture}
\qquad
\begin{tikzpicture}[scale=0.8]
  \draw[omDef, thick] plot[domain=-1:1, samples=30]
    ({\x}, {0.55*\x*\x});
  \draw[omThm, thick] plot[domain=-1:1, samples=30]
    ({0.45*\x}, {-0.55*\x*\x + 0.18*\x});
  \node[below, font=\small] at (0,-1.15)
    {अतिपरवलयीय परवलयज};
\end{tikzpicture}

{\small सप्ताहांत समस्या में वर्गीकृत नौ \omterm{pb:b2:surfaces:1}{द्विघातज} पृष्ठों में से
चार, अपनी छायाकृतियों तथा एक स्तर वक्र (लाल) से रेखांकित: परिबद्ध
दीर्घवृत्तज, अपनी कमर वाला द्विगुण रेखित एकपृष्ठी अतिपरवलयज,
दीर्घवृत्तीय परवलयज का कटोरा, और काठी, जिसके दोनों परवलयी परिच्छेद
विपरीत दिशाओं में मुड़ते हैं।}
\end{omfigure}

\begin{remark}[यह कहाँ काम आता है]
क्षेत्रफल-अवयव $\norm{\sigma_u\wedge\sigma_v}\,\dd u\,\dd v$ \cref{ch:b2:multint} का पृष्ठ-समाकल माप है, जहाँ वह ग्रीन
के सूत्र से मिलता है; \omterm{def:b2:surfaces:fff}{प्रथम आधारभूत रूप} द्विघात रूपों के किसी क्षेत्र
का प्रतिरूप है, जिसका \omterm{def:b2:funcseq:def}{बिंदुवार} अध्ययन \cref{ch:b2:quadratic} के औज़ारों से होता है;
और इस अध्याय की सप्ताहांत समस्या वर्णक्रमीय प्रमेय से सारे \omterm{pb:b2:surfaces:1}{द्विघातज}
पृष्ठ वर्गीकृत कर देती है। तृतीय वर्ष का खंड \omterm{def:b2:diffcalc:differential}{अवकल} रूपों तथा अपसरण
प्रमेय के साथ पृष्ठों पर लौटता है, और आंतरिक \omterm{thm:b2:curves:frenet2d}{वक्रता} --- मुड़ाव के
विषय में $E, F, G$ जो जानते हैं --- असली अवकल ज्यामिति का प्रवेशद्वार है।
\end{remark}

\section{अभ्यास}

\begin{exercise}[$\star$]\label{exo:b2:surfaces:1}
दिखाइए कि शंकु $z = \sqrt{x^2 + y^2}$ (उसका शीर्ष हटाकर) के सारे \omterm{def:b2:surfaces:tangent}{स्पर्श समतल} शीर्ष
से होकर जाते हैं। \emph{($\sigma(\theta, r) = (r\cos\theta, r\sin\theta, r)$, $r > 0$ से प्राचलन कीजिए।)}
\end{exercise}

\begin{exercise}[$\star$]\label{exo:b2:surfaces:2}
दीर्घवृत्तज $\frac{x^2}{a^2} +
\frac{y^2}{b^2} + \frac{z^2}{c^2} = 1$ का, उसके किसी बिंदु $(x_0, y_0,
z_0)$ पर, \omterm{def:b2:surfaces:tangent}{स्पर्श समतल}
ढूँढ़िए।
\end{exercise}

\begin{exercise}[$\star$]\label{exo:b2:surfaces:3}
हेलिकॉइड $\sigma(u, v) = (v\cos u,\
v\sin u,\ au)$, $a > 0$, के लिए $E, F, G$ परिकलित कीजिए, और टुकड़े $0 \leq u \leq
2\pi$,
$0 \leq v \leq 1$, का \omterm{def:b2:surfaces:area}{क्षेत्रफल} किसी समाकल के रूप में ($\int\sqrt{v^2 + a^2}\,\dd v = \frac12\bigl(v\sqrt{v^2+a^2} +
a^2\ln(v + \sqrt{v^2 + a^2})\bigr) + C$ का उपयोग करके उसका
मूल्यांकन कीजिए)।
\end{exercise}

\begin{exercise}[$\star\star$]\label{exo:b2:surfaces:4}
(टोरस) $xz$-समतल में केंद्र $(R, 0, 0)$ तथा त्रिज्या $r < R$ वाले वृत्त को
$z$-अक्ष के परितः घुमाकर मिलने वाले टोरस का प्राचलन कीजिए:
\[
\sigma(\theta, \psi) = \bigl((R + r\cos\psi)\cos\theta,\
(R + r\cos\psi)\sin\theta,\ r\sin\psi\bigr).
\]
$E, F, G$ परिकलित कीजिए, नियमितता जाँचिए, और दिखाइए कि \omterm{def:b2:surfaces:area}{क्षेत्रफल}
$4\pi^2 R r$ है (पाप्पुस: माध्य परिधि $2\pi R$ गुणा वृत्त की \omterm{def:b2:curves:length}{लंबाई} $2\pi r$)।
\end{exercise}

\begin{exercise}[$\star\star$]\label{exo:b2:surfaces:5}
दिखाइए कि $f \in \mathcal{C}^1(K)$ के ग्राफ़ का \omterm{def:b2:surfaces:area}{क्षेत्रफल} $\iint_K \sqrt{1 + f_x^2 + f_y^2}\,\dd x\,\dd y$ है, और उसे चक्रिका
$x^2 + y^2 \leq 1$ के ऊपर परवलयज के टुकड़े $z = \frac12(x^2 + y^2)$ के लिए परिकलित कीजिए (ध्रुवीय
निर्देशांक, \cref{ch:b2:multint})।
\end{exercise}

\begin{exercise}[$\star\star$]\label{exo:b2:surfaces:6}
त्रिज्या $R$ के गोले (गोलीय चार्ट) पर खींचे गए किसी वक्र
$\gamma(t) = \sigma(u(t), v(t))$ के लिए $u = \theta(t)$, $v =
\varphi(t)$। उसकी \omterm{def:b2:curves:length}{लंबाई} $\theta, \varphi$ में किसी समाकल के रूप में
लिखिए और सिद्ध कीजिए कि एक ही याम्योत्तर $\theta = \theta_0$ के दो बिंदुओं को
जोड़ने वाले वक्रों में याम्योत्तर चाप सबसे छोटा है।
\emph{(समाकल्य को नीचे से $R\abs{\varphi'}$ द्वारा परिबद्ध कीजिए।)}
\end{exercise}

\begin{exercise}[$\star\star\star$]\label{exo:b2:surfaces:7}
(गोले की अभिलंब रेखाएँ) मान लीजिए $S$ कोई \omterm{def:b2:metric:connected}{संबद्ध} नियमित स्तर
पृष्ठ $\{f = c\}$ है जिसकी सारी अभिलंब रेखाएँ किसी नियत बिंदु $\Omega$ से
होकर जाती हैं। दिखाइए कि $S$ $\Omega$ पर केंद्रित किसी गोले में
समाहित है। \emph{(दिखाइए कि $S$ पर खींचे गए हर वक्र के अनुदिश $\norm{M - \Omega}^2$
का अवकलज शून्य है।)}
\end{exercise}

\begin{exercise}[$\star\star\star$]\label{exo:b2:surfaces:8}
(\omterm{def:b2:surfaces:area}{क्षेत्रफल} ज्यामितीय है) मान लीजिए $\Phi \colon U' \to U$ $\R^2$ तथा $\tilde\sigma =
\sigma \circ \Phi$ के \omterm{def:b2:metric:topology}{विवृत}
समुच्चयों के बीच कोई $\mathcal{C}^1$ अवकल समाकृतिकता है। दिखाइए कि
\[
\norm{\tilde\sigma_{u'} \wedge \tilde\sigma_{v'}}
= \abs{\det J_\Phi}\,
\norm{(\sigma_u \wedge \sigma_v)\circ\Phi} ,
\]
और \cref{ch:b2:multint} के चर-परिवर्तन सूत्र का उपयोग करते हुए निष्कर्ष निकालिए कि
\cref{def:b2:surfaces:area} का \omterm{def:b2:surfaces:area}{क्षेत्रफल} चुने गए नियमित प्राचलन पर निर्भर नहीं करता।
\end{exercise}

\begin{exercise}[$\star$]\label{exo:b2:surfaces:9}
(आर्किमिडीज़ की टोप-डिब्बा प्रमेय) त्रिज्या $R$ के गोले पर
$z_1 \leq z \leq z_2$ ($-R \leq z_1 < z_2 \leq R$) वाले अक्षांशों के बीच के \emph{कटिबंध} का \omterm{def:b2:surfaces:area}{क्षेत्रफल}
$2\pi R\,(z_2 - z_1)$ है: इसे गोलीय चार्ट से सिद्ध कीजिए, और निष्कर्ष निकालिए कि
किसी कटिबंध का \omterm{def:b2:surfaces:area}{क्षेत्रफल} केवल उसकी ऊँचाई पर निर्भर करता है ---
किसी संतरे को बराबर मोटाई की फाँकों में काटने पर छिलका बराबर मात्रा
में मिलता है।
\end{exercise}

\begin{exercise}[$\star\star$]\label{exo:b2:surfaces:10}
दिखाइए कि घूर्णन पृष्ठ
$\sigma(\theta, z) = (r(z)\cos\theta,\ r(z)\sin\theta,\ z)$
($r > 0$ वर्ग $\mathcal C^1$ का) की हर अभिलंब रेखा घूर्णन-अक्ष से मिलती है, और
प्रतिच्छेद बिंदु ढूँढ़िए।
\end{exercise}

\begin{exercise}[$\star\star$]\label{exo:b2:surfaces:11}
(बेलन को खोलना) इकाई बेलन के चार्ट $\sigma(u, v) = (\cos u,\
\sin u,\ v)$ के लिए $E = G = 1$, $F = 0$: इसे
सत्यापित कीजिए, और समझाइए कि खींचे गए हर वक्र $t \mapsto
\sigma(u(t), v(t))$ की \omterm{def:b2:curves:length}{लंबाई} समतल
वक्र $t
\mapsto (u(t), v(t))$ की \omterm{def:b2:curves:length}{लंबाई} के बराबर क्यों है। निष्कर्ष निकालिए कि $(1, 0, 0)$ से
$(1, 0, 2\pi c)$ तक एक फेरा लगाने वाली कुंडली की \omterm{def:b2:curves:length}{लंबाई} $2\pi\sqrt{1 + c^2}$ है, और यह कि उन्हीं
अंत्यबिंदुओं तथा एक पूरे फेरे वाला कोई भी खींचा गया वक्र इससे छोटा
नहीं है।
\end{exercise}

\begin{exercise}[$\star\star\star$]\label{exo:b2:surfaces:12}
मान लीजिए $S = \{f = c\}$ कोई \omterm{def:b2:metric:compact}{संहत} नियमित स्तर पृष्ठ है और $M_0 \in S$ मूल बिंदु से
अधिकतम दूरी वाला कोई बिंदु। दिखाइए कि $\nabla f(M_0)$ $\vect{OM_0}$ के संरेख है ---
अर्थात् सबसे दूर वाले बिंदु पर अभिलंब त्रिज्य है। इसे दीर्घवृत्तज
$\frac{x^2}{a^2} + \frac{y^2}{b^2} + \frac{z^2}{c^2} = 1$ ($a >
b > c > 0$) पर लगाइए: वे सारे बिंदु ढूँढ़िए जहाँ अभिलंब त्रिज्य
है, और सबसे दूर वालों को पहचानिए।
\end{exercise}

\section{समस्या: $\R^3$ के द्विघातजों का
वर्गीकरण}

\begin{problem}[{सप्ताहांत समस्या --- हर द्विघातज पृष्ठ,
वर्णक्रमीय प्रमेय से छाँटा हुआ}]\label{pb:b2:surfaces:1}
\emph{द्विघातज}\index{द्विघातज} $\R^3$ में किसी द्वितीय-घात बहुपद
\[
q(X) = X^{\mathsf T}\!AX + 2\,\langle b, X\rangle + c,
\qquad A \in \mathcal S_3(\R),\ A \neq 0,\ b \in \R^3,\
c \in \R .
\]
का शून्य समुच्चय है। इस अध्याय के चित्रों के पृष्ठ --- गोले,
दीर्घवृत्तज, काठियाँ, शंकु, बेलन --- सब \omterm{pb:b2:surfaces:1}{द्विघातज} हैं। यह समस्या
उन \emph{सबका} वर्गीकरण करती है: वर्णक्रमीय प्रमेय
(\cref{thm:b2:quadratic:spectral}) द्विघात भाग को सीधा कर देती है, \omterm{def:b2:affine:subspace}{आनत} स्थानांतरण
(\cref{ch:b2:affine}) रैखिक भाग को सोख लेते हैं, और जो बचता है वह मानक रूपों
की एक छोटी, पूरी सूची है।

\medskip
\noindent\textbf{भाग I --- समानयन यंत्र।}
\begin{enumerate}
  \item मान लीजिए $X = PY + t$, जहाँ $P \in O(3)$ और $t \in \R^3$ (कोई दृढ़
        निर्देशांक-परिवर्तन)। दिखाइए कि $q(PY + t) =
        Y^{\mathsf T}\!A'Y + 2\langle b', Y\rangle + c'$, जहाँ
        \[
        A' = P^{\mathsf T}\!AP, \qquad
        b' = P^{\mathsf T}(At + b), \qquad
        c' = q(t) .
        \]
        निष्कर्ष निकालिए कि $A$ का \omterm{def:b2:reduction:eigen}{वर्णक्रम} (अतः उसकी कोटि
        और चिह्नांक) समीकरण का दृढ़ निश्चर है, और समझाइए कि
        किसी दिए गए \omterm{pb:b2:surfaces:1}{द्विघातज} का समीकरण केवल किसी अशून्य अदिश
        गुणक तक ही क्यों निर्धारित होता है।
  \item वर्णक्रमीय प्रमेय का उपयोग करते हुए दिखाइए कि किसी
        घूर्णन के बाद समीकरण $\sum_i \lambda_i y_i^2 +
        2\sum_i\beta_iy_i + c = 0$ बन जाता है, जहाँ $\lambda_1, \lambda_2,
        \lambda_3$
        $A$ के \omterm{def:b2:reduction:eigen}{अभिलक्षणिक मान} हैं।
  \item $\lambda_i \neq 0$ वाले हर $i$ के लिए $\beta_iy_i$ को किसी स्थानांतरण
        ($y_i \mapsto y_i -
        \beta_i/\lambda_i$) से सोख लीजिए। $\operatorname{rank} A = r$ होने पर समानीत समीकरण
        लिखिए: $\sum_{i\leq r}\lambda_i
        z_i^2 + 2\sum_{i > r}\beta_iz_i + c'' = 0$।
  \item समीकरण $q = 0$ वाले \omterm{pb:b2:surfaces:1}{द्विघातज} का \emph{केंद्र} ऐसा बिंदु
        $\Omega$ है कि सब $X$ के लिए $q(2\Omega - X) = q(X)$: $\Omega$ में बिंदु
        परावर्तन समीकरण को, अतः पृष्ठ को, सुरक्षित रखता है।
        दिखाइए कि $q(2\Omega - X)
        - q(X) = -4\langle A\Omega + b, X\rangle + 4\langle
        A\Omega + b, \Omega\rangle$, और निष्कर्ष निकालिए: केंद्र ठीक
        $A\Omega = -b$ के हल हैं; वे विद्यमान हैं तभी जब $b \in \operatorname{im}A$, और
        केंद्र अद्वितीय है तभी जब $A$ प्रतिलोमनीय हो।
\end{enumerate}

\medskip
\noindent\textbf{भाग II --- केंद्रीय \omterm{pb:b2:surfaces:1}{द्विघातज}
($\operatorname{rank}A = 3$)।} यहाँ समानीत समीकरण $\lambda_1z_1^2 + \lambda_2z_2^2 + \lambda_3z_3^2 = \delta$ है।
\begin{enumerate}[resume]
  \item ज़रूरत पड़ने पर $-1$ से गुणा करके मान लीजिए कि कम से कम
        दो $\lambda_i > 0$ हैं। संभावनाएँ गिनिए: $\delta > 0$, $= 0$, $< 0$ के साथ
        चिह्नांक $(3, 0)$, तथा $\delta > 0$, $= 0$, $<
        0$ के साथ चिह्नांक $(2, 1)$;
        परिणामी छह समुच्चयों को नाम दीजिए (दीर्घवृत्तज, बिंदु,
        रिक्त समुच्चय, एकपृष्ठी अतिपरवलयज, शंकु, द्विपृष्ठी
        अतिपरवलयज) और हर एक को उसके यूक्लिडीय मानक रूप में
        ($\frac{x^2}{a^2} + \frac{y^2}{b^2} +
        \frac{z^2}{c^2} = 1$, इत्यादि) रखिए।
  \item $x^2 + y^2 + z^2 + 4xy + 4yz + 4zx = 1$ का वर्गीकरण कीजिए: दिखाइए कि $A = 2J - I$, जहाँ $J$
        सब-एक आव्यूह है, \omterm{def:b2:reduction:eigen}{वर्णक्रम} $\{5, -1, -1\}$ परिकलित कीजिए, और
        अक्ष $\R(1,1,1)$ के परितः किसी घूर्णी द्विपृष्ठी अतिपरवलयज
        को पहचानिए।
  \item (रेखक) एकपृष्ठी अतिपरवलयज $\frac{x^2}{a^2} + \frac{y^2}{b^2} - \frac{z^2}{c^2}
        = 1$ के लिए
        \[
        \Bigl(\frac xa - \frac zc\Bigr)
        \Bigl(\frac xa + \frac zc\Bigr)
        = \Bigl(1 - \frac yb\Bigr)\Bigl(1 + \frac yb\Bigr)
        \]
        का गुणनखंडन कीजिए और पृष्ठ पर पड़ी सरल रेखाओं के दो
        एकप्राचलीय कुल उत्पन्न कीजिए।
  \item दिखाइए कि एकपृष्ठी अतिपरवलयज के हर बिंदु से होकर हर
        कुल की ठीक एक रेखा जाती है: अर्थात् पृष्ठ \emph{द्विगुण
        रेखित} है।
  \item एकपृष्ठी अतिपरवलयज का \emph{अनंतस्पर्शी शंकु} $C : \frac{x^2}{a^2} + \frac{y^2}{b^2}
        - \frac{z^2}{c^2} = 0$
        है। $\rho^2 = \frac{x^2}{a^2}
        + \frac{y^2}{b^2}$ के साथ दिखाइए कि अतिपरवलयज का कोई भी बिंदु
        $C$ से अधिक से अधिक $c\bigl(\rho -
        \sqrt{\rho^2 - 1}\bigr) = \frac{c}{\rho +
        \sqrt{\rho^2 - 1}}$ दूरी पर है, अतः पृष्ठ
        अनंत पर अपने शंकु से चिपक जाता है। समतलों $x = \pm a$ से
        अतिपरवलयज के परिच्छेद क्या हैं?
\end{enumerate}

\medskip
\noindent\textbf{भाग III --- कोटि $2$ और कोटि $1$:
परवलयज, बेलन, समतल।}
\begin{enumerate}[resume]
  \item मान लीजिए $\operatorname{rank}A = 2$, कहिए $\lambda_1,
        \lambda_2 \neq 0 = \lambda_3$। प्रश्न 3 से आरंभ करके
        $\beta_3 \neq 0$ (प्रश्न 4 से, कोई केंद्र नहीं) अथवा $\beta_3 = 0$
        (केंद्रों की एक रेखा) के अनुसार दो स्थितियों में
        बाँटिए, और
        \[
        \lambda_1z_1^2 + \lambda_2z_2^2 + 2\beta_3z_3 = 0
        \qquad\text{या}\qquad
        \lambda_1z_1^2 + \lambda_2z_2^2 + c'' = 0 :
        \]
        तक समानयन कीजिए --- पहली स्थिति में दीर्घवृत्तीय या
        अतिपरवलयीय परवलयज, दूसरी में केंद्रीय शांकवों के ऊपर
        बेलन (अथवा प्रतिच्छेदी समतलों का युग्म, कोई रेखा,
        रिक्त समुच्चय)।
  \item दिखाइए कि काठी $z = xy$ अतिपरवलयीय परवलयज है: $xy$-समतल
        में $\pi/4$ से घुमाकर $z = \tfrac12(u^2 - v^2)$ तक पहुँचिए, जो मापन तक \cref{fig:b2:surfaces:saddle}
        का पृष्ठ है।
  \item दिखाइए कि काठी $z = xy$ दो रेखा-कुल $\{x = x_0,\ z = x_0y\}$ तथा $\{y = y_0,\ z
        = xy_0\}$ वहन
        करती है, और हर बिंदु से होकर हर कुल की ठीक एक रेखा
        जाती है: यही दूसरा द्विगुण रेखित \omterm{pb:b2:surfaces:1}{द्विघातज} है।
  \item $\R^3$ में $x^2 + y^2 - 2x + 4y + 3 = 0$ का वर्गीकरण कीजिए (वर्ग पूरे कीजिए;
        किसी लंब वृत्तीय बेलन को पहचानिए, और उसकी अक्ष तथा
        त्रिज्या बताइए)।
  \item अब मान लीजिए $\operatorname{rank}A = 1$, कहिए $\lambda_1
        \neq 0 = \lambda_2 = \lambda_3$। केंद्रक समतल के भीतर
        घुमाकर और स्थानांतरण करके
        \[
        \lambda_1z_1^2 + 2\beta z_2 = 0 \quad (\beta \neq 0)
        \qquad\text{या}\qquad
        \lambda_1z_1^2 + c'' = 0 :
        \]
        तक समानयन कीजिए --- कोई परवलयी बेलन, अथवा समांतर समतलों
        का युग्म, कोई द्विगुण समतल, या रिक्त समुच्चय। $(x + y)^2 =
        z$ का
        पूरा वर्गीकरण कीजिए (मानक रूप, स्थानांतरण-अपरिवर्त्यता
        की अक्ष)।
\end{enumerate}

\medskip
\noindent\textbf{भाग IV --- वर्गीकरण प्रमेय।}
\begin{enumerate}[resume]
  \item भाग I से III को एक प्रमेय में जोड़िए: \emph{$\R^3$ का हर
        \omterm{pb:b2:surfaces:1}{द्विघातज} किसी दृढ़ गति से ठीक एक मानक रूप पर भेजा जाता
        है}। सत्रह \omterm{def:b2:affine:subspace}{आनत} प्रकार सूचीबद्ध कीजिए (रिक्त रूपांतर
        तथा अपभ्रष्ट समुच्चय गिनते हुए), और नौ \omterm{pb:b2:surfaces:1}{द्विघातज}
        \emph{पृष्ठ} अलग से बताइए: दीर्घवृत्तज, एकपृष्ठी तथा
        द्विपृष्ठी अतिपरवलयज, शंकु, दीर्घवृत्तीय तथा
        अतिपरवलयीय परवलयज, और दीर्घवृत्तीय, अतिपरवलयीय तथा
        परवलयी बेलन।
  \item वर्गीकरण \emph{कलनविधि} लिखिए: $(A,
        b, c)$ दिया हो, तो आप कौन
        सी राशियाँ, किस क्रम में परिकलित करते हैं, और कौन सी
        शाखा कौन सा प्रकार तय करती है? न्यायसंगत ठहराइए कि हर
        पग प्रभावी है (किसी \omterm{def:b2:quadratic:adjoint}{सममित} $3\times3$ आव्यूह के \omterm{def:b2:reduction:eigen}{अभिलक्षणिक
        मान}, कोटि, $A\Omega = -b$ की हल-योग्यता)।
  \item कलनविधि को $x^2 + y^2 + z^2 - 2xy - 2yz -
        2zx = 1$ पर चलाइए: दिखाइए कि $A = 2I - J$ का \omterm{def:b2:reduction:eigen}{वर्णक्रम}
        $\{2, 2,
        -1\}$ है और निष्कर्ष निकालिए: $\R(1,1,1)$ के परितः कोई घूर्णी
        एकपृष्ठी अतिपरवलयज।
  \item उसे $x^2 + y^2 - z^2 - 2x + 4y + 2z + 4 = 0$ पर चलाइए: केंद्र ढूँढ़िए और \omterm{pb:b2:surfaces:1}{द्विघातज} पहचानिए।
  \item उसे $x^2 + 4xy + y^2 = 2z$ पर चलाइए: $xy$-खंड का विकर्णन कीजिए ($u = \frac{x+y}{\sqrt2}$,
        $v =
        \frac{x-y}{\sqrt2}$) और \omterm{pb:b2:surfaces:1}{द्विघातज} पहचानिए।
  \item यूक्लिडीय बनाम \omterm{def:b2:affine:subspace}{आनत}। दिखाइए कि मानक रूप वाले दो
        केंद्रीय \omterm{pb:b2:surfaces:1}{द्विघातज} \emph{दृढ़ता से} तुल्य हैं तभी जब
        उनकी गुणांक-सूचियाँ एक ही हों (क्रमचय तथा समीकरण पर
        किसी साझे धनात्मक अदिश तक), जबकि \emph{\omterm{def:b2:affine:subspace}{आनत} रूप से}
        केवल चिह्नांक-आँकड़े बचते हैं: हर दीर्घवृत्तज गोल गोले
        का \omterm{def:b2:affine:subspace}{आनत} प्रतिबिंब है। कौन सी प्रमेय आश्वस्त करती है कि
        रास्ते भर चिह्नांक बदल नहीं सकता (\cref{thm:b2:quadratic:sylvester})?
\end{enumerate}

\medskip
\noindent\textbf{भाग V --- लाभांश।}
\begin{enumerate}[resume]
  \item दिखाइए कि किसी \omterm{pb:b2:surfaces:1}{द्विघातज} का किसी \omterm{def:b2:affine:subspace}{आनत} समतल से हर
        परिच्छेद उस समतल का कोई शांकव (संभवतः अपभ्रष्ट) है।
        काठी $z = xy$ के, समतलों $z = c$ ($c \neq 0$ तथा $c = 0$) से, परिच्छेद
        पहचानिए।
  \item कौन से \omterm{pb:b2:surfaces:1}{द्विघातज} पृष्ठ सरल रेखाएँ समाहित करते हैं?
        दिखाइए कि दीर्घवृत्तज, द्विपृष्ठी अतिपरवलयज तथा
        दीर्घवृत्तीय परवलयज कोई नहीं समाहित करते \emph{($q$
        को किसी रेखा तक प्रतिबंधित कीजिए, और द्विपृष्ठी
        स्थिति के लिए कोशी--श्वार्ज़ असमिका का उपयोग कीजिए)};
        कि शंकु तथा बेलन एक कुल से रेखित हैं; और निष्कर्ष
        निकालिए कि द्विगुण रेखित \omterm{pb:b2:surfaces:1}{द्विघातज} पृष्ठ ठीक एकपृष्ठी
        अतिपरवलयज और अतिपरवलयीय परवलयज हैं।
  \item जब केवल \emph{\omterm{def:b2:affine:subspace}{आनत}} प्रकार चाहिए, तब \omterm{thm:b2:quadratic:gauss}{गाउस समानयन}
        (\cref{thm:b2:quadratic:gauss}) विकर्णन से सस्ता पड़ता है। प्रश्न 17 को गाउस की
        कलनविधि से फिर कीजिए और चिह्नांक $(2, 1)$ जाँचिए; गाउस
        कौन सी यूक्लिडीय सूचना खो देता है?
  \item किसी \omterm{def:b2:quadratic:adjoint}{सममित} आव्यूह के सारे \omterm{def:b2:reduction:eigen}{अभिलक्षणिक मान} वास्तविक होते
        हैं; दिखाइए कि इसके फलस्वरूप $A$ के अभिलक्षणिक मानों
        के \emph{चिह्न} देकार्त के चिह्न-नियम से \omterm{def:b2:reduction:charpoly}{अभिलक्षणिक
        बहुपद} में ही पढ़े जा सकते हैं, और इसे प्रश्न 6 पर
        सत्यापित कीजिए: $\chi_A(\lambda) = \lambda^3 - 3\lambda^2
        - 9\lambda - 5$ में ठीक एक चिह्न-परिवर्तन है, अतः
        चिह्नांक $(1, 2)$।
  \item संश्लेषण। कलनविधि का कुछ पंक्तियों में सारांश दीजिए;
        (क) वर्णक्रमीय प्रमेय, (ख) केंद्र-समीकरण $A\Omega = -b$, (ग)
        सिल्वेस्टर की जड़त्व प्रमेय, (घ) \omterm{thm:b2:quadratic:gauss}{गाउस समानयन} --- हर एक
        की ठीक-ठीक भूमिका बताइए। वही यंत्र $\R^2$ में क्या देता
        है, और $\R^n$ में क्या बदल जाता है?
\end{enumerate}
\end{problem}
